\section{Introduction}
\label{sec:intro}

Industrial time-series monitoring---detecting faults, degradations, and anomalies in equipment ranging from semiconductor process tools to wind turbines and rotating machinery---has become a standard deep-learning application. Yet the field's practical bottleneck is not fitting a detector to one machine; it is making that detector work on the \emph{next} machine, plant, or industry, where operating conditions, sensor configurations, sampling rates, and fault statistics all change. This survey is about that bottleneck: \emph{cross-domain industrial time-series anomaly detection}.

The literature has grown along four largely separate lines. Transfer-learning research adapts detectors across domains through feature alignment, discrepancy minimization, or invariance objectives \cite{dacad2024,2310_14790,2503_08749,2407_13978,2502_00545}. Deployment-reliability research supplies distribution-free \emph{thresholds} through conformal prediction and its online, adaptive descendants \cite{lei2018,aci2021,conformalpid2023,dtaci2022,2212_03463,2407_06390}. A third line studies robustness mechanisms at inference time---test-time adaptation and parameter-efficient fine-tuning \cite{petsa2025,2312_11976,2412_08435}. A fourth line asks whether generic sequence models can transfer zero-shot, with time-series foundation models and their critiques \cite{chronos2024,2502_03383,2505_11349,2510_26777}. Surveys exist for time-series anomaly detection \cite{2002_04236,2501_15196,2308_00393,2408_03747}, for graph-based variants \cite{2302_00058,2307_03759}, for concept drift \cite{2506_15831}, and for transfer learning in anomaly detection broadly \cite{survey2023}; none organizes the cross-\emph{industry} problem jointly with the calibration and evaluation layers, which is where we find the field's most consequential and least examined issues.

Three commitments distinguish this survey. First, we treat the \emph{deployed threshold}, not the anomaly score, as the object that must transfer, and we follow how the conformal-prediction line makes that object reliable under distribution shift. Second, we take evaluation seriously: a growing body of work shows that headline metrics and protocols in this field are easy to get wrong \cite{2109_05257,2308_13068,2409_13053,2405_12096,2511_18739,2607_11969}, and we collect the specific pitfalls---point adjustment, label protocols, below-chance rankings---that cross-domain studies must handle. Third, we record negative and null results as first-class findings, because they shape method design: invariance objectives and test-time adaptation frequently fail to beat simple baselines on real industrial shifts \cite{2204_01637,2212_11080,2505_11349}, and our companion benchmarking study finds that calibration and detection quality are largely decoupled across domain families. Surveys that bury such results do readers a disservice.

The remainder of the paper is organized as follows. Section~\ref{sec:setting} formalizes cross-domain anomaly detection and the shift taxonomy. Section~\ref{sec:datasets} tabulates datasets and benchmarks. Section~\ref{sec:transfer} organizes transfer mechanisms. Section~\ref{sec:calibration} surveys reliability and calibration, including the alarm-budget problem of online conformal methods. Section~\ref{sec:evaluation} analyzes evaluation practices and pitfalls. Section~\ref{sec:open} lists open problems, and Section~\ref{sec:conclusion} concludes.

\section{Problem Setting and Notation}
\label{sec:setting}

\subsection{Tasks and Domains}
Industrial monitoring spans three closely related tasks on the same signals: unsupervised anomaly detection (point/range labels on raw sensor streams), fault diagnosis (class labels for known fault modes), and remaining-useful-life (RUL) prediction (a regression target along run-to-failure trajectories). Cross-domain settings study a model trained on one or more \emph{source domains} and deployed on a \emph{target domain} without (or with few) target labels \cite{survey2023,dacad2024}. A domain here is defined by the physical asset, operating condition, sensor set, or industry family: the field's hardest and most under-evaluated setting is leave-one-\emph{industry}-out transfer (e.g., fab to wind to machinery), as opposed to the more common leave-one-condition-out setup \cite{2503_23060,2502_00545}.

\subsection{What Changes Across Domains}
We find it useful to separate four change axes, because different transfer mechanisms address different axes:
\begin{itemize}
\item \textbf{Spectral shift}: the dominant frequency content of the signal changes---different machines rotate at different speeds, and datasets are sampled at rates from 1\,Hz (process sensors) to 128\,kHz (bearing accelerometers). Spectral shift is a leading driver of foundation-model failure on industrial series \cite{freqmatters2025,foundational2025}, and naive resampling erases the very signal a cross-domain evaluation should measure.
\item \textbf{Covariate/conditional shift}: sensor sets, scaling, and normal operating distributions differ; feature-space adaptation targets this axis \cite{dacad2024,2310_14790}.
\item \textbf{Label/prior shift}: anomaly rates and fault taxonomies differ; anomaly-injection and contrastive methods target this axis \cite{dacad2024,2411_00914}.
\item \textbf{Concept drift (temporal)}: the target itself changes over time; online calibration and drift-aware detection target this axis \cite{2506_15831,2412_08435,2410_10041}.
\end{itemize}
Real cross-industry deployment changes all four axes simultaneously, which is why single-axis methods (e.g., discrepancy minimization) underperform when evaluated across industries rather than across conditions \cite{2503_23060}.

\subsection{A Pipeline View}
We frame the field as a pipeline with four components that the literature optimizes unevenly: (i) a \emph{score producer} (reconstruction, forecasting residual, contrastive discriminator, foundation model); (ii) a \emph{transfer mechanism} (feature alignment, invariance objectives, adapters, experts); (iii) a \emph{calibration layer} that maps scores to decisions under a nominal error level; and (iv) an \emph{evaluation protocol} joining datasets, labels, and metrics. Most papers invest in (i)--(ii); Sections~\ref{sec:calibration}--\ref{sec:evaluation} argue that (iii)--(iv) determine whether the resulting system is deployable.
